\section{Training Observability：Loss Spike 只是症状}

训练不是提交 job 后等待数周。研究假设可能被 data corruption、optimizer instability、precision overflow、code regression、network degradation 或 checkpoint mismatch 破坏。单看 global loss 无法诊断，因此需要 model、optimizer、data 与 system 四个同步 telemetry plane。

\subsection{四类 Telemetry 与 anomaly triage}

\term{loss spike} 是训练 loss 突然异常上升，可能短暂恢复，也可能预示 run 已失效。Model plane 记录 loss、gradient/activation statistics；optimizer plane 记录 learning rate、moment/update norm；data plane 记录 batch source、token mixture 与 corruption；system plane 记录 utilization、network、storage 和 hardware error。

\lecturefigure{06-training-observability.png}{训练诊断需要 model、optimizer、data 与 system 四类同步 telemetry。}{本地字幕 17:20--20:10；概念重绘。}

\begin{knowledgebox}{读图：Correlation 先于 root cause}
当 loss spike 出现，先把异常时间对齐 data batch、optimizer update、collective retry 和 device error，再提出 hypothesis。若只看 loss，团队可能错误降低 learning rate；若只看 hardware log，又可能忽略特定数据触发的数值不稳定。所有 telemetry 必须使用一致 run/step ID，并保存 code/data/config version。
\end{knowledgebox}

\begin{warningbox}{Dashboard 绿色不代表实验有效}
GPU utilization、throughput 和 job status 可以全部正常，而训练目标已因错误 mask、污染数据、label shift 或 optimizer state 损坏而失真。Frontier program 需要 scientific invariants：held-out loss、gradient range、data distribution、checkpoint checksum 与 periodic evaluation。
\end{warningbox}

\subsection{Checkpoint、replay 与 resume}

\term{checkpoint} 保存恢复训练所需的参数、optimizer state、scheduler state、random state 与 data cursor；它不同于 activation checkpointing，后者是在单个 forward/backward 内用重算节省显存。\term{deterministic replay} 尽可能用相同 code、batch 和 state 重现异常，以判断是 transient infrastructure fault 还是 recipe defect。

\lecturefigure{07-checkpoint-recovery.png}{训练恢复要保留科学含义：detect、checkpoint、replay、diagnose，再 resume 或修复。}{本地字幕 17:20--19:10；概念重绘。}

\begin{knowledgebox}{读图：恢复点必须包含 data cursor}
只保存 model weights 会丢失 optimizer moments、LR schedule、randomness 和数据位置，恢复后曲线可能悄悄改变。Replay 若在同一 batch 再次 spike，优先检查 data/recipe；若换硬件后消失，可能是 transient fault；若无法重现，则要保留不确定性并提高监控，而不是宣布问题已解决。
\end{knowledgebox}

\teachervoice{Mann 用“像看护病人一样看训练”说明 on-call 心态：loss curve 的异常可能在数小时后才显现，回滚又会损失昂贵进度。课堂里的重点不是戏剧化值班，而是训练本身已经成为需要 SRE discipline 的生产系统。}

\subsection{本章小结}

Training observability 把模型实验与系统运行连接起来。Loss spike 是入口，四类 telemetry、完整 checkpoint 与 controlled replay 才能把 anomaly 变成可修复知识。

